% SPDX-License-Identifier: CC-BY-4.0
% Copyright (c) 2026 Martin Schröder <info@swedishembedded.com>

\section{Limitations}\label{sec:limitations}

The structure claims show that definitions, constants, tests and messages exist at the pinned
commits; they do not show that the code behaves accordingly. Behaviour is established only by
the experiments, which are pilots: one run each per setting, one generator family, small
corpora, 30 facts per generation, one policy size per run. The training settings of E2 were
chosen on two earlier pilots with other seeds, not on the registered run's own measurement, but
they were chosen by hand and no search over them was made.

No run released a candidate. In every run the improvement, retention and serve checks passed or,
in one run, the serve check could not start its server, and the anchor check refused; so the
alias move, the serving manifest of a real release, a second generation that continues a
champion with replay, the retention check against an earlier release and the rollback to a first
release are described from the code (\cref{sec:gate}) and are not demonstrated end to end here.
The serve check was measured on refused candidates only. The anchor bound allows a loss of less
than one of 40 questions, so a single changed answer decides it; we did not test a larger anchor
suite or other training settings that might keep it, beyond the two settings and two policy sizes
reported. The facts are numbers, which may be why they leak into numeric answers elsewhere; a
domain of non-numeric facts is not measured.

The held-out families of a run are facts the policy is never taught, about a seventh to a quarter of
the facts, and the teacher's answer was not verified for a further share of them
(\cref{sec:defects}); the facts measured in the probes include both, so the counts of facts answered
are lower bounds of what the trained records taught. Probes are graded by numbers appearing in the
answer; a different grader was not compared. The experiments ran on binaries built from the commits
of \cref{sec:repro}, which differ between the first and second wave and, for the repaired-verifier
rerun, from the registered run's; each result names its wave. The training qualification of the 27B
model used random weights for the long sequences and a handful of examples for the real checkpoint;
it shows that steps fit and run, not that the adapter changes generation, and no sequence longer
than 16384 tokens was run.

The cause of the generator's missing subjects is localised, not resolved: the rule is not at fault
beyond the markup case, the generator at two sizes names a subject its question omits, and a changed
brief was not tested. Corpus expansion beyond the first pass (variants, retention queue,
rehearsal, claims) is described from the code; variants and rehearsal were exercised as stage
counts and as a post hoc setting, and none of the mechanisms was evaluated for the value it adds.

\input{generated/limitations}
